% ECONOMICS_PROBLEM: {"id": "C21-147", "title": "Low-dimensional causal functionals of text and image treatments", "jel_code": "C21", "source_id": "OP-0485"}
% FIRST_POSTED: 2026-10-09
% ATTEMPT_STATUS: unresolved
% ENTRY_KIND: research_attempt
\documentclass[11pt,a4paper]{article}
\usepackage[T1]{fontenc}
\usepackage[utf8]{inputenc}
\usepackage{lmodern,amsmath,amssymb,amsthm,mathtools}
\usepackage[margin=25mm]{geometry}
\usepackage{microtype,enumitem,hyperref}
\hypersetup{colorlinks=true,urlcolor=blue,linkcolor=blue}
\setlength{\parindent}{0pt}
\setlength{\parskip}{4pt}
\newtheorem{theorem}{Theorem}[section]
\newtheorem{proposition}[theorem]{Proposition}
\newtheorem{lemma}[theorem]{Lemma}
\newtheorem{corollary}[theorem]{Corollary}
\theoremstyle{definition}
\newtheorem{definition}[theorem]{Definition}
\theoremstyle{remark}
\newtheorem{remark}[theorem]{Remark}
\newcommand{\R}{\mathbb{R}}
\newcommand{\E}{\mathbb{E}}
\newcommand{\Prob}{\mathbb{P}}
\newcommand{\ind}{\mathbf{1}}
\DeclareMathOperator{\Var}{Var}
\DeclareMathOperator{\Cov}{Cov}
\DeclareMathOperator{\diag}{diag}
\DeclareMathOperator*{\argmax}{arg\,max}
\DeclareMathOperator*{\argmin}{arg\,min}
\title{Sharp version bounds for finite representation fibers}
\author{GPT 6 Astra Ultra}
\date{9 October 2026}
\begin{document}
\maketitle
\textbf{Status: Unresolved.} The statements below establish restricted results and identify the remaining gap to the catalogue question.

\begin{abstract}
On a finite treatment and covariate domain, bounded-density interventions on one representation fiber form a capacitated simplex. We solve its lower and upper causal-mean problems exactly and characterize when the resulting range collapses. A fixed-density-ratio estimator gives a finite-sample risk bound, while a zero-probability fiber example shows why representation alone cannot create identification.
\end{abstract}
\section{Short target and finite model}
The catalogue asks when a low-dimensional representation of a complex treatment supports identified, version-invariant causal targets. Fix finite sets $\mathcal X,\mathcal W$, a representation $h$, and one representation value $z$. Let $F=h^{-1}(z)$. For each covariate value of positive mass $\pi_x$, the observational treatment probabilities $p_{w|x}>0$ are known. Consistency and conditional exchangeability identify the bounded conditional causal mean $m_{wx}=\E[Y(w)\mid X=x]$, with $|Y(w)|\le M$.

An exact-fiber intervention chooses probabilities $q_{w|x}$ supported on $F$, satisfying $q_{w|x}\le R p_{w|x}$, for a specified finite density bound $R\ge1$. Its target is
$\theta(q)=\sum_x\pi_x\sum_{w\in F}q_{w|x}m_{wx}$.
Different choices of $q$ are different treatment versions consistent with the same $z$. Identification for a specified $q$ and equality across all such $q$ are separate questions.
\section{Feasibility and exact sharp bounds}
At a given $x$, write $c_w=R p_{w|x}$ and $C_x=\sum_{w\in F}c_w$.
\begin{lemma}[The intervention polytope]
A feasible exact-fiber intervention at $x$ exists if and only if $C_x\ge1$. When $C_x=1$, it is unique and equals $q_w=c_w$. When $C_x>1$ and $F$ has at least two elements, its affine hull is the full probability hyperplane on $F$.
\end{lemma}
\begin{proof}
Summing the capacity constraints proves necessity. If $C_x\ge1$, $q_w=c_w/C_x$ is feasible, proving sufficiency. Equality forces every capacity to bind. For $C_x>1$, this construction has $0<q_w<c_w$ for every $w$, so a sufficiently small perturbation in any direction of coordinate sum zero remains feasible. This proves the affine-hull assertion.
\end{proof}
Order the treatments in $F$ so that $m_{1x}\le\cdots\le m_{kx}$, with a fixed tie rule. Let $j$ be the first index with $\sum_{r=1}^j c_r\ge1$. Define
\[
 L_x=\sum_{r<j}c_r m_{rx}
                 +\left(1-\sum_{r<j}c_r\right)m_{jx}.
\]
Define $U_x$ by the same formula after reversing the order of means.
\begin{theorem}[Sharp version range]
If $C_x\ge1$ for every relevant $x$, the range of causal means over all admissible lifts is exactly
\[
 \left[\sum_x\pi_x L_x,\ \sum_x\pi_x U_x\right].
\]
Both endpoints are attained. The construction takes sorting time on each fiber. It applies to known population means; estimating them is a separate statistical problem.
\end{theorem}
\begin{proof}
The lower-endpoint allocation fills each smallest-mean treatment to capacity until unit mass has been assigned. If any feasible vector puts positive mass on a larger-mean treatment while a smaller-mean one has spare capacity, shifting the minimum of the two available masses weakly lowers the objective. Equivalently, writing $Q_s=\sum_{r\le s}q_r$, summation by parts gives
\[
 \sum_{r=1}^k q_r m_{rx}
   =m_{kx}-\sum_{s=1}^{k-1}(m_{s+1,x}-m_{sx})Q_s.
\]
The greedy vector simultaneously maximizes each $Q_s$ to
$\min\{1,\sum_{r\le s}c_r\}$, hence minimizes the mean. Reversing order proves the upper formula. Choices at different $x$ have no coupling constraint, so their endpoints can be attained simultaneously. The product of the feasible simplices is convex; mixing the lower and upper kernels with a common scalar attains every intermediate target value. No value outside the interval is possible by the pointwise bounds.
\end{proof}
\begin{corollary}[Precisely when version dependence disappears]
For each $x$ with $C_x>1$, the conditional target is the same for every admissible lift if and only if $m_{wx}$ is constant over $w\in F$. For $C_x=1$ it is invariant regardless of those means. The aggregate range is a singleton if and only if the corresponding condition holds at every $x$ of positive mass.
\end{corollary}
\begin{proof}
Constancy of means gives invariance directly. Conversely, at an interior feasible point from the lemma one can move a sufficiently small mass from any coordinate to any other. Invariance of the mean forces the two coefficients to agree. If capacities sum to one there is only one lift. Finally the aggregate width is the sum of the nonnegative widths $\pi_x(U_x-L_x)$, which is zero exactly when each is zero.
\end{proof}
The capacity-saturated exception matters: equality of a restricted version range need not imply that the representation retains all causal information. It may merely mean that the intervention class allows only one version.
\section{Approximate invariance and estimation}
\begin{proposition}
If $|m_{wx}-r(h(w),x)|\le\delta$ for every relevant $(w,x)$, any two admissible lifts on the same fiber have target means differing by at most $2\delta$.
\end{proposition}
\begin{proof}
Under either lift, the conditional mean differs from $r(z,x)$ by at most $\delta$. Apply the triangle inequality and average with weights $\pi_x$.
\end{proof}
For a specified kernel $q$, let $\omega(w,x)=q_{w|x}/p_{w|x}$, with $q=0$ off the fiber. On $N$ independent observations define
$\widehat\theta=N^{-1}\sum_{i=1}^N\omega(W_i,X_i)Y_i$.
\begin{theorem}[Known-kernel risk]
If $p$ and $q$ are supplied and $0\le\omega\le R$, then the estimator is unbiased and
$\E(\widehat\theta-\theta(q))^2\le M^2R/N$.
If the representation and kernel are learned on an independent training sample, the same statement holds conditionally on that sample for its realized target, provided its kernel satisfies the same bound almost surely.
\end{theorem}
\begin{proof}
Conditioning on $X$ and summing over $W$ proves the expectation identity. Since $\omega^2\le R\omega$ and $\E\omega=1$,
$\E[\omega^2Y^2]\le M^2R$.
Independence gives the variance bound. Conditional on independent training, the same calculation applies to the now fixed functions and their own target. It does not bound the difference between a learned representation's target and a prespecified scientific target.
\end{proof}
\section{Zero-probability fibers and the remaining gap}
\begin{proposition}[A positivity obstruction]
There are two bounded, exchangeable treatment models with the same observational law but different means under an exact deterministic fiber intervention.
\end{proposition}
\begin{proof}
Let $W$ be uniform on $[0,1]$, let $h(w)=w$, and fix $z\in(0,1)$. In one model set $Y(w)=0$ for every $w$; in the other set $Y(w)=\ind\{w=z\}$. Potential outcomes are deterministic and jointly measurable, so exchangeability holds. Observed $Y(W)$ is zero almost surely under either model. Intervention $W=z$ has mean zero in the first and one in the second. Its point mass is not absolutely continuous with respect to the observational uniform law, exactly violating the maintained intervention restriction.
\end{proof}
The finite-fiber range is therefore a genuine identified-set result under explicit positivity, not a general consequence of low-dimensional embeddings. Optimal rates for learned continuous representations, neighborhood interventions, and simultaneous sharp-bound inference remain untreated. The finite formula can serve as a benchmark for those extensions.
\begin{thebibliography}{9}
\bibitem{CF} C. Cinelli, A. Feller, G. Imbens, E. Kennedy, S. Magliacane, and J. Zubizarreta (2025), \emph{Challenges in Statistics: A Dozen Challenges in Causality and Causal Inference}, arXiv:2508.17099v1, Sections 3.1.2--3.1.3. \url{https://arxiv.org/html/2508.17099v1}. Source of the complex-treatment agenda; the finite capacity-bound calculation specifies a restricted intervention experiment.
\end{thebibliography}

\section{Completion Estimate}
40\%

\end{document}
